\section{A tame horizontal logarithmic comparison}\label{sec:tame}

We prove a first-derivative comparison in a tame character direction. Its
weak form shows that the derivative vanishes when the projected Heegner
trace is torsion, without assuming analytic rank one. Its rank-one form
then identifies the logarithm of the
Kato class with the classical leading coefficient, with its exact rational
normalization.

\subsection{The comparison and its horizontal data}

Let \(E/\Q\) be the non-CM curve on the branch under consideration, with
primitive newform \(f\), conductor \(N\), rational differential \(\omega\),
and absolute primitive real Betti-cycle period \(\Omega_0\). Thus \(E\)
may be the chosen modular quotient or its indicated positive twist.
The CM adaptations will be made in the CM comparison. Fix a support
\(S\) containing \(2,\infty\), the bad primes, and the fixed orders and
smoothing data. Write \(S_f\) for its finite part. We use the Kato system
\eqref{eq:K1}, \eqref{eq:K4} at symbol evaluation \(1\), at this support
and after adjoining one prime; supplementary primes can be removed by norm.
Let \(z\) be its base smoothed class in ordinary cohomology, and put
\begin{equation}\label{tame:d0}
 d_0=\pm\Delta(0)\prod_{q\in S_f}L_q(f,1)^{-1}\ne0.
\end{equation}
The smoothing and omitted Euler factors in this expression are the actual
ones on the branch. All signs below come from a consistent choice of
orientation. We identify \(T=T_2E\) and \(T^\vee(1)\) by the principal
polarization, and use \(\log_\omega\) on finite classes on either side.

\begin{theorem}[Tame logarithmic comparison]\label{tame:logarithm}
Suppose \(\an(E)=1\), and let \(P\in E(\Q)\) be nontorsion. Then
\[
 \frac{\log_\omega z}{(\log_\omega P)^2}
   =\pm d_0\,\frac{L'(f,1)}{\Omega_0H(P)}.
 \tag{T1}\label{eq:T1}
\]
The complex-side quotient is rational, and is viewed in a fixed dyadic
embedding. Here \(H\) is the height fixed in the introduction.
\end{theorem}

The logarithm-squared relation in \eqref{eq:T1} belongs to the framework
proposed by Perrin-Riou~\cite[Section~3.3]{PerrinRiou1993}.
Bertolini--Darmon--Venerucci prove such a relation, up to a nonzero
rational factor, at odd semistable primes
\cite[Theorem~A]{BertoliniDarmonVenerucci2022}. The tame comparison
below determines the stated dyadic normalization.

For the auxiliary construction we require only that \(w(E)=-1\).
Choose an imaginary quadratic field \(K\), of odd fundamental discriminant
\(-D<-4\), split at \(2N\) and all further required fixed primes, such that
\[
 L(f\otimes\varepsilon,1)\ne0,\qquad \varepsilon=\varepsilon_K.
\]
The local prescriptions and nonvanishing are those of the classical
twist theorem recalled above \cite{FH95}. Put \(O=\mathcal O_K\).
Choose an integral polynomial \(\mathcal A\) in finitely many \(T_p\),
with \(p\nmid2ND\), which has eigenvalue \(\mathcal A_f\ne0\) at \(f\)
and kills every other eigensystem at level \(N\). Include among the systems
to be killed the Eisenstein types unramified away from \(N\), of conductor
dividing \(N\), and the weight-two trivial pair. Multiplicity one and
rational characteristic or minimal polynomials give such a polynomial
after clearing denominators.

Choose also \(\ell\nmid6ND\), away from the operators in \(\mathcal A\),
inert in \(K\), with \(a_f(\ell)\ne0\). This is possible by
Chebotarev and the open-image theorem: the inert condition specifies
an open coset in the Tate image, on which trace is not identically zero
\cite{Serre72}. All coefficient tests below use a fixed finite collection
of indices prime to \(ND\), with odd \(\ell\)-valuation, and powers of
one moving prime \(r=r_i\).

\begin{definition}[Horizontal data]\label{tame:data}
A horizontal sequence consists of primes \(r_i\), integers \(m_i\to\infty\),
and even surjective residue exponents
\[
 \lambda_i:(\Z/r_i)^\times\longrightarrow\Z/2^{m_i},
\]
with the following properties. The primes avoid the fixed data,
\(r_i\equiv1\pmod{ND}\), \(r_i\to1\) dyadically, and \(r_i\) splits in
the Hilbert class field of \(K\) and in any fixed twisting field in use.
The characters kill every prime of \(S_f\) and of the fixed smoothing
integers \(c,d\), including \(2\). We may require them also to kill any
specified finite list, in particular \(3\), the primes of \(D\), the
tested indices and Hecke operators, and the bad primes of the fixed
models and lifts below. Chosen arithmetic Frobenius lifts converge in
the Tate representation to a determinant-one element with trace
\(a_*\ne\pm2\) and with two eigenvalues that are not roots of unity.

As a global character, \(\lambda_i\) has value \(\lambda_i(q)\) on a
uniformizer at \(q\ne r_i\), is zero at infinity, and is the negative
residue character at \(r_i\), with uniformizer value zero. Characters
over extensions are obtained by norm.
\end{definition}

We use the simultaneous Artin limits of the arithmetic diagram
conventions, and write \(t=1+u\) in this section. At a character evaluation,
\(\chi(x)=t^{\lambda_i(x)}\). Let \(z_i\) be the Kato classes with
coefficients in \(\Z_2[t]/(t^{2^{m_i}}-1)\), on the Tate dual with the
scalar action, or its inverse, as prescribed by orientation. Their
local coefficients at \(2\) are split, since \(\lambda_i(2)=0\).
Let \(Z_i(t)\) denote the polynomial of dual-exponential coordinates
relative to \(\omega\). Invert the variable if necessary so that
\begin{equation}\label{tame:reciprocity}
 \begin{split}
 z_i(1)&=e_rz,\qquad e_r=1-a_f(r)/r+1/r,\\
 Z_i(t)&=d_0\,L(f,\bar\chi,1)/\Omega_0
       \qquad(\ord\chi>2).
 \end{split}
\end{equation}
The first equality is rational in cohomology. These are the good-prime
norm relation and \eqref{eq:K4} at square order, with unnormalized weights.
There is no omitted \(r\)-factor at the displayed primitive characters.
The positive quadratic realization, if present, uses the same relations
because \(r\) splits in that fixed field. The sign from
\eqref{eq:K3}--\eqref{eq:K5} is common at each stage; fix it in \(d_0\)
in the limit. Ordinary local coordinates apply coefficient by coefficient,
so these classes and scalar polynomials have bounded denominators.
We have \(Z_i(1)=0\), and denote the limiting linear coefficient by \(Z'\).

Let \(P_X\in J_0(N)(K)\) be the single Hilbert-class trace based at
\(\infty\), with the oriented ideal \(\mathfrak n\) of norm \(N\) used in
\eqref{eq:GZ}. Fix a rational modular parametrization
\(\phi:X_0(N)\to E\), of degree \(\delta\), and put
\[
 Q_X=\phi_*P_X,\qquad
 C=\frac{\sqrt D}{8\pi^2(f,f)_N}.
\]
The Petersson integral here is unnormalized.

\begin{proposition}[Spectral-height comparison]\label{tame:spectral}
For the horizontal data of Definition~\ref{tame:data}, and the tame
height defined below, one has
\begin{equation}\label{tame:spectral-identity}
 C\Omega_0L(f\otimes\varepsilon,1)\frac{Z'}{d_0}
   =\pm(2-a_*)\,\frac{h_\lambda^K(Q_X,Q_X)}{\delta}.
\end{equation}
This assertion requires the odd functional sign and the nonzero
companion value, but does not assume \(\an(E)=1\) or that \(E(\Q)\)
spans the rational Selmer group. In particular, if \(Q_X\) is torsion,
then \(Z'=0\).
\end{proposition}

We construct the directions and heights first. The analytic and
intersection calculations then prove Proposition~\ref{tame:spectral}.
Only afterward will we use analytic rank one to deduce
Theorem~\ref{tame:logarithm}.

\subsection{Directions, biextensions, and cup reciprocity}

We first isolate the choice of a character direction from its height
interpretation. The group-theoretic input is a nonzero cohomology class
and a primitive for its polarized self-cup. This also permits the same
choice when the class is not known to come from a rational point.

\begin{lemma}[A direction detecting a self-cup primitive]
\label{tame:cochain-direction}
Let \(V=T_2E\otimes\Q_2\), and let
\(x\in Z^1(G_{\Q,S},V)\) represent a nonzero class. Let \(\beta\)
be a cocycle on the Tate dual representing \(\operatorname{pol}([x])\),
and suppose that a continuous
cochain \(k\in C^1(G_{\Q,S},\Q_2(1))\) satisfies
\[
                         dk=-\beta\cup x.
\]
Fix a finite normal extension \(F/\Q\) enforcing the splitting and
congruence conditions of Definition~\ref{tame:data}, including the
Hilbert class field of \(K\), the fixed twisting data, and \(E[2]\).
Let \(\mathcal P\) be a finite list of rational primes containing the
fixed support to be killed, and put
\[
 L=F\bigl(\mu_{2^\infty},\,q^{1/2^n}:q\in\mathcal P,\ n\ge1\bigr).
\]
There is \(g\in G_L\) whose Tate action \(g_T\) has determinant one,
trace \(a_*\ne\pm2\), and non-root-of-unity eigenvalues, such that
\begin{equation}\label{tame:eta}
 \eta=k(g)-\beta(g)g_T(g_T-1)^{-1}x(g)\ne0.
\end{equation}
The horizontal data can be chosen with Frobenius lifts converging to
\(g\) on these cochains and all fixed data, and with residue characters
vanishing on \(\mathcal P\).
\end{lemma}

\begin{proof}
After rescaling \(x,\beta\) by one common dyadic integer and \(k\)
by its square, compactness makes all three cochains integral. Their
fixed-support condition makes them unramified at the moving primes.

\emph{The Tate image and the additive evaluations.}
The quotient defining \(L/F\) has bounded derived length. The
open-image theorem and iterated closed commutators therefore leave
a deep subgroup of \(\mathrm{SL}_2(\Z_2)\) in the image of \(G_L\):
the corresponding iterated Lie brackets contain \(\mathfrak{sl}_2\).

Restriction of \([x]\) to \(G_F\) is nonzero, since restriction followed
by corestriction is multiplication by \([F:\Q]\). Its restriction to
\(G_{F(T)}\) is still nonzero: otherwise it would inflate from the open
Tate image, whose first cohomology on \(V\) vanishes by a central
homothety. On \(G_{F(T)}\) its values are additive and equivariant, so
irreducibility makes their span the full Tate plane.

The same holds on \(G_{L(T)}\). Indeed, if the restriction there were
zero, those additive evaluations would factor through
\(\operatorname{Gal}(L(T)/F(T))\). Since \(F(T)\) contains all dyadic
roots of unity, this radical quotient has cyclotomic conjugation.
A deep special-linear subgroup consequently acts trivially on that
quotient, whereas it has no trivial quotient on the standard plane.
This contradicts the full span just obtained. Irreducibility again
makes the restricted span the full plane.

\emph{The scalar detected by the direction.}
Choose \(g\in G_L\) with the asserted Tate action, for example in a
deep split torus. On the joint Tate kernel coboundaries vanish, so
\(\beta\) agrees with \(\operatorname{pol}(x)\). The cochain identity gives,
up to the fixed sign,
\[
 k(h_1h_2)=k(h_1)+k(h_2)+\langle x(h_1),x(h_2)\rangle.
\]
A commutator there has \(x=\beta=0\) and \(k\)-value
\(\pm2\langle x(h_1),x(h_2)\rangle\). The full span supplies a
nonzero such value. Multiplying \(g\) by this commutator changes
\(\eta\) without changing its Tate action or its \(x,\beta\) values.
This gives \eqref{tame:eta}. Only nonvanishing is used; the factor
\(2\) is retained.

\emph{Frobenius and residue characters.}
Chebotarev gives primes \(r_i\) and Frobenius lifts \(g_i\) approaching
\(g\) on all the fixed data, with roots and radicals fixed to depths
sufficient for \(m_i\). Choose \(r_i\equiv1\pmod{2^{m_i+1}}\).
In compatible primitive-root bases define
\begin{equation}\label{tame:residue-exponent}
 a^{(r_i-1)/2^{m_i}}\equiv
 \zeta_{2^{m_i}}^{\lambda_i(a)}\pmod{r_i}
 \qquad(r_i\nmid a).
\end{equation}
These residue exponents are surjective and define even characters.
Fixing the prescribed radicals makes them vanish on \(\mathcal P\),
and the Frobenius conditions give all the remaining horizontal data.
\end{proof}

\begin{lemma}[A direction with nonzero height]\label{tame:direction}
If \(P\in E(\Q)\) is nontorsion, the horizontal data can be chosen so
that \(h_\lambda^\Q(P,P)\ne0\).
\end{lemma}

\begin{proof}
Write \(y=\operatorname{pol}(P)\). Let \(G_y\) be the Barsotti--Weil
extension of \(E\) by \(\mathbb G_m\) defined by the Poincaré biextension
at \(y\), and choose \(B\in G_y(\Q)\) above \(P\). Such a lift exists
because its fiber is a line torsor over \(\Q\).
Choose a splitting of the underlying modules
\(T_2G_y=\Z_2(1)\oplus T\); its Galois matrix has off-diagonal block
\(\beta(g)g_T\), where \(\beta\in Z^1(T^\vee(1))\).
This is, up to the fixed sign, the Kummer class of \(y\). To see the
identification integrally, at order \(b\) the torsor of division points
\(by'=y\) maps to the torsor of splittings by pushing out
\(G_{y'}[b]\) along the \(b\)-power map on the kernel. Changing \(y'\)
by \(b\)-torsion changes the splitting by its Weil pairing. These
identifications commute with reduction.

The Kummer cocycle of \(B\) therefore has coordinates \((k,x)\), with
\(x\) the Kummer class of \(P\), and
\begin{equation}\label{tame:biextension-cocycle}
 dk=-\beta\cup x.
\end{equation}
The data and representatives are unramified outside fixed support,
by good integral models and the étaleness of division by \(2\) away
from that support. The class of \(x\) is nonzero because \(P\) is
nontorsion. Lemma~\ref{tame:cochain-direction} supplies horizontal
data with \(\eta\ne0\). The height construction in the next lemma
identifies \(h_\lambda^\Q(P,P)=\pm\eta\), proving the assertion.
\end{proof}

\begin{lemma}[Tame heights and trace compatibility]\label{tame:heights}
Let \(F'\supset K\) contain the ring class field \(H_{r^j}\), \(j\ge1\).
The norm of \(\lambda_i\) to \(F'\) is finite-unramified. On the
polarized abelian varieties in use it defines a bilinear
biextension height \(h_i^{F'}\) modulo \(2^{m_i}\).
After multiplying the first argument by one fixed power of \(2\),
there are compatible heights on \(E/K\) and \(E/\Q\), even at the
ramified character place \(r_i\). Their rational limits satisfy
trace compatibility and
\[
 h_\lambda^K(P_1,P_2)=2h_\lambda^\Q(P_1,P_2)
 \qquad(P_1,P_2\in E(\Q)).
\]
For the data in Lemma~\ref{tame:direction},
\(h_\lambda^\Q(P,P)=\pm\eta\).
\end{lemma}

\begin{proof}
At either split place of \(r\), a local unit whose image dies in the
ring class quotient must be a diagonal congruence modulo \(r^j\),
up to the global units \(\{\pm1\}\). This follows directly from the
idele description for the order \(\Z+r^jO\). Local Artin reciprocity
and norm show that the even character becomes unramified over
\(H_{r^j}\), and hence over \(F'\). It was already unramified elsewhere.
Bad reduction and fixed model primes have character zero.

Choose any rational lift in the Poincaré biextension at the two
arguments, with polarization on the second. Sum the local character
valuations of this lift relative to the rigidified integral fibers
at good places. Reciprocity makes the sum independent of the lift.
The two integral biextension laws give bilinearity and adjoint
compatibilities. On a relative Jacobian, degree-zero horizontal
divisors with disjoint generic support compute this valuation,
up to the common sign, by their intersection lengths. Indeed the
Poincaré line under self-duality is the Deligne pairing, or its
inverse according to orientation, and the change of its symbol
lattice is exactly the intersection. Adding a whole vertical fiber
has zero pairing with a degree-zero divisor. Only models outside
the fixed killed support are needed.

For the downstairs heights, observe that
\[
 \#E(\F_{r_i})=r_i+1-a_f(r_i)\longrightarrow2-a_*\ne0.
\]
Its two-adic valuation is bounded. Choose \(2^b\) which removes all
these finite two-primary parts, so the scaled first argument \(X\)
lies in the uniquely \(2\)-divisible factor of the local points.
For a lift \(\mathcal B\) at \((X,Y)\), use at \(r_i\) the character
of the scalar
\[
 \mathcal B/([2^{m_i}]_1\mathcal B'),
\]
where \(\mathcal B'\) is any lift at
\(([2^{m_i}]^{-1}X,Y)\), and division is in that factor.
The biextension law makes this independent of the choice and
bilinear modulo \(2^{m_i}\). If the normed character becomes
unramified in an extension, compare with integral lifts to recover
the valuation definition there. Division in the first argument
can be performed downstairs; norming the lifts and scalar ratios,
using the second biextension law, proves compatibility with trace
in the second argument. At other places this is integral trace
compatibility. Taking limits and dividing by \(2^b\) defines the
rational heights. The field degree gives the stated factor \(2\)
on rational arguments.

For the final assertion choose the fixed lift \(B\) above \(P\),
scaled in the first direction. It is integral away from killed
places, so only \(r_i\) contributes. Put
\[
 w_i=(g_{i,T}-1)^{-1}x(g_i).
\]
Unramifiedness gives \(dw_i=x\) locally; enlarge \(b\) so
\(2^bw_i\) is integral for every \(i\). Subtracting the coboundary
of \((0,2^bw_i)\) from the Kummer cocycle of \(B\) produces invariant
division points of the scaled first argument, in its uniquely
divisible factor. Its remaining fiber coordinate is the unramified
cocycle \(2^b(k-\beta\cup w_i)\). More concretely, a corrected
division point differs from any rational lift above the divided
base point by a root of the scalar ratio defining the height.
Its Frobenius value in the root basis of
\eqref{tame:residue-exponent} is precisely the residue character,
up to the fixed sign; the valuation part of unramified Kummer
is divisible by the root order. Passing to the limit and dividing
by \(2^b\) gives \(h_\lambda^\Q(P,P)=\pm\eta\).
\end{proof}

\begin{lemma}[Cup reciprocity in analytic rank one]\label{tame:cup}
If \(\an(E)=1\) and \(P\in E(\Q)\) is nontorsion, then
\[
 Z'\log_\omega P
  =\pm(2-a_*)\frac{\log_\omega z}{\log_\omega P}
       h_\lambda^\Q(P,P).
 \tag{T2}\label{eq:T2}
\]
\end{lemma}

\begin{proof}
In analytic rank one, the classical theorem gives rational rank one
and finite \(\Sha\). The base dual exponential of \(z\) is zero by
\eqref{eq:K4}. At odd finite places the rational \(H^1\) of the
Tate coefficients is zero, by vanishing of invariants and dual
invariants and the local Euler characteristic. Thus \(z\) lies in
the rational Selmer line. Also \(\log_\omega P\ne0\), since the
kernel of the local logarithm consists of torsion.

Work modulo \(2^{m_i}\) and \(u^2\), and write a cocycle for \(z_i\)
as \(z_0+uz_1\). Then
\[
 dz_1=-\lambda_i\cup z_0.
\]
Use the opposite sign throughout for the inverse scalar action.
After one fixed common multiplier, the rational norm relation and
Selmer line allow a coboundary adjustment with
\[
 z_0=\alpha_i\beta,\qquad
 \alpha_i=\pm e_r\,\frac{\log_\omega z}{\log_\omega P}.
\]
The integral \(H^1\)-torsion in this adjustment has bounded exponent
by global torsion finiteness. Suppressing that common multiplier,
the root-coefficient two-cocycle
\[
 z_1\cup x+\alpha_i\lambda_i\cup k
\]
is closed by \eqref{tame:biextension-cocycle}.
Its local invariant at \(2\) tends to \(Z'\log_\omega P\), by
logarithm--dual-exponential local duality. Here localization has
split ordinary group coefficients before differentiation, so all
pairings are made on ordinary lifts with only fixed denominators.
At the other fixed places the character is zero and the residual
contribution has bounded torsion order; away from the support and
\(r_i\), the cup is unramified.

At \(r_i\), the trivialization \(dw_i=x\) replaces this cocycle in
cohomology by
\(\alpha_i\lambda_i\cup(k-\beta\cup w_i)\).
Local cup reciprocity identifies its invariant, up to sign, with
\(\alpha_i\) times the fiber character computed in
Lemma~\ref{tame:heights}. Global reciprocity, followed by the limit
and \(e_r\to2-a_*\), proves \eqref{eq:T2}.
\end{proof}

\subsection{The holomorphic kernel}

The height and cup constructions are now in place. To prove
Proposition~\ref{tame:spectral}, we compare the two sides of its identity
through Fourier coefficients of a weight-two modular form. Its
projection to the \(f\)-component will give the product involving
\(Z'\); its character derivative will give a weighted intersection of
Heegner divisors. The restriction to indices with odd valuation at the
inert prime \(\ell\) will make those divisors disjoint.

Choose a real primitive period \(\Omega_{\rm tw}\) for
\(E^{-D}\), with a rational differential. We will compare complex
identities in a dyadic embedding using
\begin{equation}\label{tame:period-rationality}
 C\Omega_0\Omega_{\rm tw}\in\Q^\times,\qquad
 L(f\otimes\varepsilon,1)/\Omega_{\rm tw}\in\Q^\times.
\end{equation}
The second assertion is modular-symbol rationality. For the first,
the complex integral of \(|\omega\wedge\bar\omega|\) is a rational
multiple of \(\Omega_0\) times the primitive absolute anti-invariant
period. If \(\phi^*\omega=c_\phi f\,dq/q\), then this integral is
\(8\pi^2c_\phi^2(f,f)_N/\delta\), with \(c_\phi\in\Q^\times\).
The negative-twist period is a rational multiple of the anti-invariant
period divided by \(\sqrt D\). These facts give
\eqref{tame:period-rationality}.

Choose integral ideals \(C_0\), prime to \(D\), representing
\(\Pic(O)\). Put \(a_C=\mathrm NC_0\) and
\(Q_C(x)=\mathrm Nx/a_C\), and define
\begin{equation}\label{tame:theta}
 \theta(z)=\frac12\sum_{C_0}\sum_{x\in C_0}e(Q_C(x)z),
 \qquad e(w)=\exp(2\pi i w).
\end{equation}
This weight-one form on \(\Gamma_0(D)\) has character \(\varepsilon\)
and positive coefficients
\(\rho(n)=\sum_{d\mid n}\varepsilon(d)\).
The factor \(1/2\) removes the two generators of a principal ideal;
our restriction \(-D<-4\) ensures there are no additional units.
For \(\ord\chi>2\), let \(\theta_\chi\) be its Fourier twist, of
character \(\varepsilon\chi^2\), and set
\begin{equation}\label{tame:kernel}
 \begin{split}
 E_{\chi,s}(z)
   &=\frac12\sum_{\substack{(c_1,d_1)\ne(0,0)\\ND\mid c_1}}
      \varepsilon(d_1)\bar\chi^2(c_1)
      \frac{\operatorname{Im}(z)^s}
      {(c_1z+d_1)|c_1z+d_1|^{2s}},\\
 I_\chi
   &=\left.\frac{\sqrt D}{2\pi}
       \operatorname{Tr}_{NDr^2/Nr^2}
          (\theta_\chi E_{\chi,s})\right|_{s=0}.
 \end{split}
\end{equation}
The row letters \(c_1,d_1\) here have no connection to the smoothing
integers. The trace is the unnormalized modular-form trace.

\begin{lemma}\label{tame:holomorphic}
The continuation in \eqref{tame:kernel} at \(s=0\) is a holomorphic
weight-two form of level \(Nr^2\), with trivial character.
\end{lemma}

\begin{proof}
Begin with the row sum in its domain of absolute convergence.
Row transformation gives the character opposite to that of
\(\theta_\chi\). At every integral cusp transform, the row weights
are periodic and odd under simultaneous change of sign. The
\(c_1=0\) sum is regular at \(s=0\) by oddness and Dirichlet
continuation. For \(c_1>0\), combine the two signs and apply Poisson
summation to the second variable.

The zero frequency is a periodic Dirichlet sum at \(2s\), times
\(y^{-s}\) and an integral factor regular at zero; the odd part of
that integral cancels. At a nonzero frequency the relevant transform
is that of
\[
 (u_0+ic_1y)^{-1}|u_0+ic_1y|^{-2s}.
\]
After rescaling, shift to the lines
\(\operatorname{Im}(u_0/(c_1y))=\pm1/2\) and integrate by parts.
The differentiated tails are absolutely convergent and give
exponential decay in \(c_1y\) times the absolute frequency, uniformly
near \(s=0\). At zero only positive frequencies survive, by the
pole in the lower half-plane. The resulting constants and positive
modes are holomorphic. This reasoning also justifies continuation
against cusp tests, as required in the unfolding below.
\end{proof}

\subsection{The partial trace at the discriminant primes}

We will express the coefficient of \(q^{nr_i^h}\) in \(I_\chi\)
as the group polynomial \(D_{i,h}(n)\) in the following lemma.
We compute these coefficients at
\[
 m=nr^h,\qquad h\ge1,\quad r\nmid n,\quad(n,ND)=1,\quad
 v_\ell(n)\ \text{odd}.
\]
The finite list of \(n\)'s includes those generated by applying
\(\mathcal A\) at \(n=\ell\). For a rational prime \(p\mid D\), put
\[
 \epsilon_p(a)=(a/p),\qquad
 \tau_p=\sum_{a\bmod p}\epsilon_p(a)e(a/p).
\]
Thus \(\epsilon_p\) is a residue symbol, distinguished below from
the local norm character \(\varepsilon_{K,p}\).

\begin{lemma}[The traced coefficient]\label{tame:coefficient}
At the indicated character evaluation, the coefficient of \(q^{nr^h}\)
in \(I_\chi\) is
\begin{equation}\label{tame:coefficient-formula}
 \begin{split}
 D_{i,h}(n)
  &=\frac12\sum_{\substack{C_0,\ y_0\in C_0,\ b>0\\
                 a=\mathrm Ny_0/a_C>0\\a+Nb=nr^hD}}
       \ \sum_{\substack{kl=b\\\gcd(k,l,D)=1}}
         W(k)\chi(a/D)\bar\chi^2(Nk),\\
 W(k)
  &=\prod_{\substack{p\mid D\\p\mid k}}\epsilon_p(l)
     \prod_{\substack{p\mid D\\p\nmid k}}\epsilon_p(-a_CNk).
 \end{split}
\end{equation}
Only terms with \(a\) a unit at \(r\) occur. Its linear coefficient
is the same masked sum with the character factor replaced by
\(-2\lambda_i(k)\).
\end{lemma}

\begin{proof}
\emph{The local theta transform.}
The dual of the lattice \((C_0,Q_C)\) is
\((1/\sqrt{-D})C_0\). For \(p\mid D\), write \(d=D/p\).
The \(p\)-primary discriminant subgroup is represented by
\(x=dy/\sqrt{-D}\), \(y\in C_0\). Since \(C_0\) is prime to \(p\),
conjugation is the identity on its residue field, and
\(\mathrm Ny\equiv t^2\pmod p\) for the residue \(t\) of \(y\).
Its quadratic form is therefore
\[
 Q_C(x)=\frac{d\,\mathrm Ny}{pa_C},\qquad
 q_p t^2/p,\quad q_p\equiv d/a_C\pmod p.
\]
The square class of \(q_p\) is that of \((a_Cd)^{-1}\).

Let \(S_0=\left(\begin{smallmatrix}0&-1\\1&0\end{smallmatrix}\right)\)
and \(T_0=\left(\begin{smallmatrix}1&1\\0&1\end{smallmatrix}\right)\).
Choose \(a\equiv-1\pmod p\) and
\(a\equiv0\pmod{(D/p)Nr^2}\). The word
\(T_0^aS_0T_0^aS_0^{-1}T_0^a\) is \(S_0\) modulo \(p\) and identity
on the other components. The finite theta transform has
\[
 T_{x,y}=\delta_{x,y}\zeta^{q_px^2},\qquad
 F_s(x,y)=p^{-1/2}\zeta^{s2q_pxy},\qquad
 \zeta=e(1/p),\ s=\pm1.
\]
The scalar in the \(S_0\)-operator cancels against its inverse.
Thus the word on this component is
\(T^{-1}F_sT^{-1}F_{-s}T^{-1}\), whose \((x,z)\)-entry is
\[
 \frac1p\sum_y
 \zeta^{-q_p(x^2+y^2+z^2)+s2q_py(x-z)}
 =\frac1p\sum_y\zeta^{-q_py^2}\zeta^{-2q_pxz}.
\]
Completing the square proves the equality. In particular its zero
row and zero column are constant. Counting squares in \(\F_p\)
gives the exact scalar
\begin{equation}\label{tame:kappa}
 \kappa_p=\frac1p\sum_y e(-q_py^2/p)
          =p^{-1}\tau_p\epsilon_p(-a_CD/p).
\end{equation}
The other primary components are unchanged. Hence the partial
\(S_0\)-branch replaces the local lattice by
\(\mathfrak p^{-1}C_0\), with scalar \(\kappa_p\).

\emph{Trace representatives and character weights.}
For a subset \(J\) of the primes dividing \(D\), put
\(D_J=\prod_{p\in J}p\), \(D_I=D/D_J\).
The partial words just constructed give a representative
\(\gamma_J\in\Gamma_0(Nr^2)\) which is \(S_0\) at the primes of
\(J\) and identity at the others. Its branch of the trace has
representatives \(\gamma_JT_0^j\), where \(j\) runs through the
residues modulo \(D_J\) and is zero modulo \(D_I\). The product
has Fourier denominators dividing \(D_J\). Summing these translations
therefore kills nonintegral indices and multiplies an integral
index \(m\) by \(D_J\). This is one factor \(p\) for each partial
branch.

To track the Fourier twist, write
\(\gamma=\left(\begin{smallmatrix}A&B\\C_\gamma&D_\gamma\end{smallmatrix}\right)\),
where \(r^2\mid C_\gamma\), and
\(n(x/r)=\left(\begin{smallmatrix}1&x/r\\0&1\end{smallmatrix}\right)\).
Direct multiplication gives
\[
 n(x/r)\gamma=\gamma'n(z'/r),\qquad
 z'\equiv xD_\gamma A^{-1}\equiv xA^{-2}\pmod r.
\]
Taking \(x,z'\) divisible by \(ND\) makes \(\gamma'\) integral and
\(\gamma'\equiv\gamma\pmod D\). Thus
\(\gamma'\gamma^{-1}\in\Gamma(D)\) acts trivially on the untwisted
theta form. Expressing \(\theta_\chi\) by translates weighted by
\(\bar\chi(x)/\tau(\bar\chi)\), the coefficient of index
\(U\in D^{-1}\Z\) is consequently multiplied by
\[
 \tau(\bar\chi)^{-1}\sum_{x\bmod r}\bar\chi(x)e(Uz'/r)
   =\chi(U)\chi(A^{-2}).
\]
The representatives divisible by \(D\) make this finite sum
unambiguous for the fractional index. Characters are extended by zero.

For the Eisenstein rows, old and new coordinates satisfy
\((c_0,d_0)=(c_1,d_1)\gamma^{-1}\), so
\(c_0\equiv c_1A^{-1}\pmod r\).
The factor \(\bar\chi^2(c_0)=\bar\chi^2(c_1)\chi^2(A)\)
cancels the extra theta factor. We retain
\(\chi(U)\bar\chi^2(c_1)\).
At an identity prime \(p\mid D\), the row has \(p\mid c_1\)
and weight \(\epsilon_p(d_1)\). At a partial \(S_0\)-prime it has
\(p\mid d_1\) and weight \(\epsilon_p(c_1)\). Also \(N\mid c_1\).

\emph{Poisson summation and its scalar.}
For \(c_1>0\), a positive frequency has
\(V=c_1l/D\), \(l>0\). Poisson summation at \(s=0\) gives
\(-2\pi i/D\) times the positive finite Fourier transform.
Indeed the transform of \((c_1z+u_0)^{-1}\), closed clockwise
in the lower half-plane, is \(-2\pi i\,e(c_1lz/D)\).
CRT evaluates the finite transform as
\[
 \prod_{\substack{p\mid D\\p\notin J}}\tau_p\epsilon_p((D/p)l)
       \prod_{p\in J}\epsilon_p(c_1).
\]
Multiplying by \(D_J\prod_{p\in J}\kappa_p\), and by
\(\sqrt D/(2\pi)\), leaves
\[
 \frac{-i}{\sqrt D}
 \left(\prod_{p\mid D}\tau_p\epsilon_p(D/p)\right)
 \prod_{\substack{p\mid D\\p\notin J}}\epsilon_p(l)
 \prod_{p\in J}\epsilon_p(-a_Cc_1).
\]
Here the local cancellation is
\[
 \frac{p\kappa_p}{\tau_p\epsilon_p(D/p)}\epsilon_p(c_1)
       =\epsilon_p(-a_Cc_1).
\]
The parenthesized product is the primitive Gauss sum
\(i\sqrt D\) of \(\varepsilon\), since \(D\equiv3\pmod4\).
Thus the common scalar is exactly one.

\emph{The lattice restrictions.}
Put \(y_0=\sqrt{-D}\,x\),
\(a=DU=\mathrm Ny_0/a_C\), \(c_1=Nk\), and \(b=kl\).
Then \(a+Nb=mD\). The zero frequency \(V=0\) contributes nothing:
the index \(m\) cannot be a theta norm because its \(\ell\)-valuation
is odd. Also \(a=0\) would imply \(N\mid mD\), impossible since
\((m,ND)=1\) and \(N>1\).
At an identity prime, \(p\mid k\); the relation forces \(p\mid a\),
which is exactly the extra theta condition
\(y_0\in\mathfrak pC_0\). The Fourier factor requires \(p\nmid l\).
At a partial prime, \(p\nmid k\), while \(l\) is unrestricted.
Thus \(J=\{p\mid D:p\nmid k\}\), the restriction is
\(\gcd(k,l,D)=1\), and the remaining scalar is \(W(k)\).
The sole factor \(1/2\) is the original theta normalization.
This proves \eqref{tame:coefficient-formula}.

Since \(r\mid m\), the relation \(a+Nb=mD\) makes \(a\) a unit
at \(r\) exactly when \(b\) is a unit. In that case \(k,l\) are
both units. Thus \(\chi(a/D)\) gives precisely the stated mask.
For a divisor pair, the ratio of the local factors of \(W(l)\)
and \(W(k)\) is \(\epsilon_p(-a_CN)\). If neither divisor is
divisible by \(p\), use the norm congruence
\(\epsilon_p(-a_CNb)=1\) to obtain the same ratio.
Hence
\[
 W(l)=-W(k),\qquad
 \varepsilon(-a_CN)=-1.
\]
The latter uses \(\varepsilon(a_C)=1\), since \(a_C\) is an ideal
norm, and \(\varepsilon(N)=1\), since the level primes split.
The derivative of the character weight is
\(\lambda_i(U)-2\lambda_i(N)-2\lambda_i(k)\).
The first two terms are independent of the divisor \(k\) and
cancel by antisymmetry. This leaves \(-2\lambda_i(k)\).
\end{proof}

Define \(B_h(n)\), also for \(h=0\), by the same sum without the
mask, replacing its character weight by
\(-2\lambda_i(k/r^{v_r(k)})\). Put
\[
 C_h(n)=B_h(n)-B_{h-2}(n)\quad(h\ge1),\qquad B_{-1}=0.
\]
These notations may also denote their limits. A prime on \(D_h\)
will mean the limiting derivative. At the working residue precision,
the unmasked sum has the decomposition
\[
 B_h(n)=\sum_{q\ne r}\lambda_i(q)B_{h,q}(n),
\]
where \(B_{h,q}(n)\) is the same sum with weight \(-2v_q(k)\).
This follows by factoring \(k/r^{v_r(k)}\) into rational primes,
and the identity passes to the limit. We will identify \(-B_{h,q}\)
with a local intersection. Thus \(D'_h\)
is supplied by the holomorphic kernel, while \(B_h\) has an
intersection interpretation. The next recurrence removes the mask
between them. The difference defining \(C_h\) will correspond to
the cyclic part \(T_{r^h}-T_{r^{h-2}}\) of the Hecke correspondence,
with \(T_{r^{-1}}=0\).

\begin{lemma}[Removal of the mask]\label{tame:mask}
For \(h\ge2\),
\[
 C_{h+1}(n)-2C_h(n)+C_{h-1}(n)
   =D'_{h+1}(n)+2D'_h(n)+D'_{h-1}(n).
 \tag{T4}\label{eq:T4}
\]
\end{lemma}

\begin{proof}
Strip the common power \(r^j\) from \(a,b\), writing
\[
 a=a^\circ r^j,\quad b=b^\circ r^j,\quad h=h_0+j,\quad
 \min(v_ra^\circ,v_rb^\circ)=0.
\]
The off-\(r\) weights do not change: \(r\equiv1\pmod D\), and
the two primes over \(r\) are principal, so their parts can be
inserted freely in \((y_0)/C_0\).
With the off-\(r\) choices fixed, there are
\(1+j+v_ra^\circ\) allocations between the two split primes in
the theta norm, and \(1+j+v_rb^\circ\) allocations in the divisor
pair \(kl=b\). Their product gives the multiplicity in \(B_h\):
\[
 (1+j+v_ra^\circ)(1+j+v_rb^\circ).
\]
If \(h_0>0\), both primitive entries are units; the resulting
\(C_h\)-multiplicities, beginning at \(h_0\), are
\(1,4,8,12,\ldots\), whereas the masked derivative has multiplicity
one only at \(h_0\). Their second difference has the coefficients
\(1,2,1\) displayed in \eqref{eq:T4}. If \(h_0=0\), the
\(C_h\)-multiplicities are affine for \(h\ge1\), and there are
no masked terms there. These contributions have zero second
difference, proving the formula.
\end{proof}

\subsection{The spectral evaluation}

We first compute the holomorphic side of the recurrence. Applying
\(\mathcal A\), evaluating at \(n=\ell\), and taking the difference
\(h\mapsto h+2\) removes every component except the oldforms attached
to \(f\). Pairing those oldforms with the kernel expresses the
result in terms of \(Z'\) and the nonzero companion value.

Write \(a_h=a_f(r^h)\), \(a^{\rm pr}=a_f(r)\), and
\[
 \gamma_h=\lim_i\big(a_f(r_i^h)-a_f(r_i^{h-2})\big)
 \qquad(h\ge1),
\]
where \(a_f(r^{-1})=0\). These limits satisfy the determinant-one
recurrence with trace \(a_*\).

\begin{lemma}[The projected holomorphic coefficient]\label{tame:spectral-coefficient}
For all sufficiently large \(h\),
\[
 \begin{split}
 \big[\mathcal A(D'_{h+2}-D'_h)\big](\ell)
   ={}&\mathcal A_fa_f(\ell)
      \frac{\gamma_{h+2}-\gamma_h}{(a_*-2)(a_*+2)}
      C\Omega_0\,\frac{Z'}{d_0}\\
     &{}\cdot(2-a_*)L(f\otimes\varepsilon,1).
 \end{split}
 \tag{T5}\label{eq:T5}
\]
\end{lemma}

\begin{proof}
\emph{Isolation of the \(f\)-block.}
Apply \(\mathcal A\) by the usual Hecke formulas on coefficients,
evaluate at \(n=\ell\), and difference by \(h\mapsto h+2\).
On \(I_\chi\) this uses \(U_r^h(U_r^2-1)\).
A cuspidal representation of conductor exponent one at \(r\)
has the signed Steinberg eigenvalue; its complementary raises
have zero eigenvalue. The displayed operator annihilates both
for large \(h\). Conductor exponent two has \(U_r=0\).
Eisenstein series with a ramified pair at \(r\) likewise have
\(U_r=0\). All remaining unwanted old types, including the
holomorphic special forms of the trivial pair, are killed by
\(\mathcal A\). Thus only the \(f\)-block remains. Good Hecke
adjoints and multiplicity one permit its coefficients to be
computed by pairings against cusp tests.

\emph{Fricke transformation and unfolding.}
At level \(Nr^2\) use
\[
 f_j=r^jf(r^jz),\qquad 0\le j\le2.
\]
Choose integers \(d_1,d_2,d_3\) with
\[
 W_r=\begin{pmatrix}rd_1&d_2/r\\rd_3&r\end{pmatrix},\qquad
 d_2,d_3\equiv0\pmod{ND},\quad d_3\equiv1\pmod r,\quad
 r^2d_1-d_2d_3=1.
\]
Such choices follow from CRT. This determinant-one normalization
of the exact-divisor matrix sends \(f_2\) to \(f\) and fixes \(f_1\),
by level-\(N\) modularity. Use the same matrix before the trace
at level \(NDr^2\). The identity
\[
 n(x/r)W_r=\gamma'n(z'/r),\qquad
 z'\equiv-x^{-1}\pmod r
\]
for unit \(x\), with \(x,z'\equiv0\pmod{ND}\) and integral
\(\gamma'\equiv1\pmod D\), gives
\[
 \theta_\chi\big|W_r
   =\frac{\tau(\chi)}{\tau(\bar\chi)}\theta_{\bar\chi}.
\]
On the Eisenstein rows, write \((v,w)\) for the old coordinates.
Clearing \(r\) gives
\[
 v'=r^2(vd_1+wd_3),\qquad
 w'=vd_2+r^2w.
\]
These are a bijection onto the rows with \(NDr^2\mid v'\):
the inverse formulas are
\[
 v=v'-d_3w',\qquad
 w=d_1w'-d_2v'/r^2.
\]
The new weights are \(\varepsilon(w')\bar\chi^2(w')\), and the
scalar is \(r^{1+2s}\).

Group the row greatest common divisors \(b_1\), which are prime
to \(Dr\). The primitive row sum has level
\(N_1Dr^2\), where \(N_1=N/\gcd(N,b_1)\).
The terms with \(N_1<N\) pair to zero by newness of \(f\).
This is also true against \(f_1\): tracing down at the \(N\)-factor
in either degeneracy gives the corresponding zero trace. The
extra levels are coprime to \(N\), so CRT separates the cosets;
conjugation by \(\operatorname{diag}(r,1)\) for \(f_1\) is integral
on those representatives and permutes the same \(N\)-coset data.

The remaining gcd factor is
\(L^{(N)}(\varepsilon\bar\chi^2,1+2s)\).
Unfold the primitive rows on \(\Gamma_0(NDr^2)\), counting the
two signs once. The \(f_1\)-pairing is zero because its Fourier
support is disjoint from that of \(\theta_{\bar\chi}\).
The \(f\)-pairing has integral factor
\(\Gamma(1+s)/(4\pi)^{1+s}\) and Dirichlet series
\[
 \sum_{j\ge1}\frac{a_f(j)\rho(j)\bar\chi(j)}{j^{1+s}}.
\]
The Hecke recurrences identify this series with the product
of the two twisted \(L\)-series divided by the row gcd factor.
At \(p\mid N\), use \(\varepsilon(p)=1\) and the degree-at-most-one
primitive local polynomial; at \(p\mid D\) the ramified twist
factor is \(1\), and at \(r\) both factors are \(1\).
Since \((f,f)_{Nr^2}=r(r+1)(f,f)_N\), evaluation at \(s=0\) gives
\begin{equation}\label{tame:old-pairings}
 \begin{split}
 \frac{(I_\chi,f_2)_{Nr^2}}{(f,f)_{Nr^2}}
    &=\frac{C}{r+1}\frac{\tau(\chi)}{\tau(\bar\chi)}
          L(f,\bar\chi,1)L(f\otimes\varepsilon,\bar\chi,1),\\
 (I_\chi,f_1)_{Nr^2}&=0.
 \end{split}
\end{equation}
In particular the factors \(\sqrt D/(2\pi)\), \(r\),
\(1/(4\pi)\), and the level index \(r(r+1)\) give exactly
\(C/(r+1)\).

\emph{The oldform Gram matrix.}
In the basis \(f_0,f_1,f_2\), the normalized Gram matrix is
\((g_{|i-j|})\), where
\[
 g_0=1,\qquad
 g_1=\frac{a^{\rm pr}}{r+1},\qquad
 rg_2=a^{\rm pr}g_1-1.
\]
These follow from norm under translation and cyclic-coset
averaging, with eigenvalues for \(T_r\) and \(T_{r^2}-1\).
At \(h\ge2\), the coefficient vector after removing the
\(\ell\)-factor is
\((a_h,ra_{h-1},r^2a_{h-2})\).
The recurrence \(a_h=a^{\rm pr}a_{h-1}-ra_{h-2}\) shows that
it lies in the span of the Gram rows for \(f_1,f_2\).
The coefficient of the second of these rows is
\[
 s_h=\frac{r^2a_{h-2}-g_1ra_{h-1}}{1-g_1^2}.
\]
The denominator has bounded dyadic valuation, since
\(a_*\ne\pm2\). Taking the limit and using the determinant-one
recurrence gives
\begin{equation}\label{tame:gram-limit}
 \lim_i\frac{s_h}{r_i+1}
   =\frac{\gamma_h}{(a_*-2)(a_*+2)}.
\end{equation}

\emph{Interpolation at the central character.}
All the factors in this old-block calculation admit polynomials
with bounded denominators. The Gauss ratio is
\(\tau(\chi)^2/r\). The Gauss polynomials have coefficients in
\(\Z_2\): Frobenius on the \(r\)-th roots fixes
them because \(\lambda_i(2)=0\). At the trivial character each
Gauss polynomial has value \(-1\), so the limiting augmentation
of the ratio is \(1\).

For \(L(f,\bar\chi,1)/\Omega_0\) use \(Z_i/d_0\).
For \(L(f\otimes\varepsilon,\bar\chi,1)/\Omega_{\rm tw}\)
use the even part of the \(\chi\)-weighted additive modular-symbol
sum at the units \(x/r\), in positive Fourier Mellin convention,
multiplied by \(\tau(\bar\chi)/r\).
Before Gauss multiplication, its nontrivial value is
\(\tau(\chi)L(f\otimes\varepsilon,\bar\chi,1)/\Omega_{\rm tw}\).
At the trivial character its value is
\[
 (a^{\rm pr}-2)L(f\otimes\varepsilon,1)/\Omega_{\rm tw}.
 \]
Indeed summing the additive twists over the nonzero residues gives
\[
 r\sum_{r\mid j}\frac{a_{f\otimes\varepsilon}(j)}j
       -L(f\otimes\varepsilon,1)
   =(a^{\rm pr}-2)L(f\otimes\varepsilon,1),
\]
by the Hecke recurrence and \(\varepsilon(r)=1\), with Mellin
continuation understood. The fixed relative period lattices and
the Manin--Drinfeld theorem bound the denominators
\cite{ManinDrinfeld}. Multiplying by the trivial Gauss value
\(-1\) and letting \(r\to1\) therefore gives augmentation
\[
 (2-a_*)L(f\otimes\varepsilon,1)/\Omega_{\rm tw}.
\]
Use \eqref{tame:period-rationality} to compare all the products
in the fixed dyadic embedding.

The character identities hold outside order at most two. In the
characteristic-zero cyclic group ring, evaluation at all characters
is injective. Multiplying the difference of the two sides by
\(t^2-1\) also kills its evaluations at the two excluded characters,
so gives a group-ring identity. The denominators of both sides are
bounded, hence this identity passes to the bounded power-series
limit. That ring is a domain, so its nonzero factor \(t^2-1\)
can now be canceled. Differentiating the
resulting identity uses \(Z_i(1)=0\), so only the central value
of the second factor is needed. Equations
\eqref{tame:old-pairings} and \eqref{tame:gram-limit}, with the
difference \(h\mapsto h+2\), now give exactly \eqref{eq:T5}.
\end{proof}

\subsection{The unmasked intersection calculation}

We now compute the local sums \(B_{h,q}(n)\) in the decomposition
of \(B_h(n)\). Their weights are \(-2v_q(k)\). The divisor calculation
below gives their values, and Proposition~\ref{tame:intersection}
identifies their negatives with the corresponding intersections of
the Hilbert-class divisors. Only \(q\nmid6NDm\) need be considered:
all remaining prime factors,
and the finitely many bad model and fixed CM-data primes, have
been killed by \(\lambda_i\). We may enlarge that fixed list
whenever a model or ideal representative is chosen.

\begin{lemma}[The local divisor sum]\label{tame:local-divisor}
Fix \(C_0,y_0,b\) in the unmasked sum, and put \(d_q=v_q(b)\).
Its contribution with weight \(-2v_q(k)\), including the factor
\(1/2\), is zero unless \(q\) is inert, \(d_q\) is odd, and
\(\varepsilon_{K,p}(-bNa_C)=1\) for every \(p\mid D\).
In that case it is
\[
 -\frac{1+d_q}{2}\,\rho(b/q)\,
       2^{\#\{p\mid D:p\mid b\}}.
 \tag{T6}\label{eq:T6}
\]
\end{lemma}

\begin{proof}
Since \(\gcd(k,l,D)=1\), a full prime power at \(p\mid D\)
lies on one side of the divisor pair. Put all these powers
initially on \(l\)'s side. Then \(W(k)=-\varepsilon(k)\).
Moving \(p^{d_p}\) to the other side changes the sign by
\[
 \epsilon_p(-bNa_C/p^{d_p})
     \prod_{\substack{s\mid D\\s\ne p}}\epsilon_s(p^{d_p}),
\]
which is the local norm sign \(\varepsilon_{K,p}(-bNa_C)\),
by the product formula. At the primes \(p\mid D\) not dividing
\(b\), that sign is already positive by the norm congruence.
Thus the allocations at the discriminant primes vanish unless
all signs are positive, and otherwise contribute
\(2^{\#\{p\mid D:p\mid b\}}\).

The remaining allocations are ordinary \(\varepsilon\)-weighted
divisors. For \(d=d_q\), set
\[
 A_q(d)=\sum_{j=0}^d\varepsilon(q)^j,\qquad
 M_q(d)=\sum_{j=0}^d(-2j)\varepsilon(q)^j.
\]
If \(q\) is split, or if \(q\) is inert and \(d\) is even, then
\(M_q(d)=-dA_q(d)\). The weighted sum is therefore proportional
to the undifferentiated sum, which is zero by \(W(l)=-W(k)\).
If \(q\) is inert and \(d\) is odd, then \(M_q(d)=d+1\);
all the other local factors give \(\rho(b/q)\), whose \(q\)-factor
is one. The initial minus sign and the factor \(1/2\) give
\eqref{eq:T6}.
\end{proof}

\begin{proposition}[Intersection realization]\label{tame:intersection}
The negative of the local sum in Lemma~\ref{tame:local-divisor},
summed over \(C_0,y_0,b\), is the intersection of the oriented
Hilbert-class sum based at \(\infty\) with \(T_m\) of the same
sum based at \(0\), at the places of \(K\) over \(q\), weighted
by their residue degrees over \(q\).
\end{proposition}

\begin{proof}
The generic supports are disjoint because \(v_\ell(m)\) is odd
and \(\ell\) is inert. The cusps are disjoint at the good model
places under consideration and retain their types under the
prime-to-\(N\) Hecke correspondences; CM points have potentially
good reduction, so there are no cusp intersections there.

At a split good prime \(q\), the endomorphism algebra of the
ordinary CM reduction is \(K\). A matching homomorphism would
be linear or conjugate-semilinear over \(K_\ell\), between Tate
lattices free of rank one over \(O_\ell\). Its degree has even
\(\ell\)-valuation, contradicting the choice of \(m\).
The intersection is consequently zero in this case.

Suppose \(q\) is inert. Write
\(\mathcal W=W(\overline{\F}_q)\).
After omitting the killed bad definition and ramification
primes of the fixed CM objects and level data, we have good
lifts over \(\mathcal W\). Classical CM describes the Hilbert
orbit, with its fixed type and orientation, by tensoring one
object with ideal representatives \(I\); the level subgroup
is cut out by \(\mathfrak n\).

\emph{Coarse intersections and lifting lengths.}
For each ordered pair of such objects, its intersection over
\(\mathcal W\) is one half the sum, over degree-\(m\)
homomorphisms of their reductions respecting level, of the
individual lifting lengths. Here is the normalization.
Since \(q\nmid m\), every subgroup of order \(m\) in the
reduction lifts uniquely, and \(T_m\) uses all these subgroups
with their multiplicities. After passage to a fine auxiliary
level, each deformation disk is smooth with tame automorphism
action, since \(q\ge5\). The coarse disk is the quotient by
this action modulo its generic kernel \(\{\pm1\}\).
Averaging an integral parameter with its tangent character
linearizes the effective action. The tangent character is
faithful: a tame kernel acting trivially on the parameter
would act trivially on the disk, contrary to the generic
automorphism group. The valuation of the coarse-parameter
difference is therefore the sum of the framed congruence
lengths over these translates. A framed congruence is exactly
a lift of the specified special isomorphism. Thus counting
all special isomorphisms gives division by \(2\).
At an inert prime the residue degree over \(\Q\) is \(2\);
this cancels that division in the weighted intersection of
the proposition.

This is the proper-intersection formula of
\cite[Theorem~5.1]{Conrad03}; its hypotheses hold here because
\((m,N)=1\) and the inert-\(\ell\) norm obstruction excludes
generic intersections. We retain the calculation to identify
its precise weight in the present divisor sum.

\emph{The quaternion lattice.}
Let \(\mathcal R\) be the supersingular maximal order of the
base object. Its quaternion algebra has the description
\[
 \mathcal B=K+Kj_0,\qquad j_0x=\bar xj_0,\qquad
 j_0^2=-c_{\rm B}<0.
\]
At a finite prime \(s\ne q\), choose a generator of the base
Tate lattice over \(O_s\). Maximality identifies \(\mathcal R_s\)
with its full integral endomorphisms, and \(j_0\) acts as
\(b_s\) times conjugation, with \(\mathrm Nb_s=-c_{\rm B}\).
At \(q\) use the unique division maximal order; \(v_q(c_{\rm B})\)
is odd. The fractional ideal \(\mathfrak c\) defined by
\[
 vb_s\in O_s\ (s\ne q),\qquad
 c_{\rm B}\mathrm Nv\in q\Z_q
\]
has norm \(q/c_{\rm B}\).

For tensor objects indexed by \(I,J_1\), the integral Hom
lattice is \(J_1\mathcal RI^{-1}\), with degree
\((\mathrm NI/\mathrm NJ_1)\) times reduced norm. The degree
factor follows from the tensor inclusion isogenies of ideals
of given norms, also after reduction. Count pairs by writing
\(J_1=C_0I\). For \(u_0+vj_0\) in rational Hom, put
\(y_0=\sqrt{-D}\,u_0\), \(z_0=\sqrt{-D}\,v\).
The lattice and level conditions become
\begin{equation}\label{tame:quaternion-lattice}
 \begin{gathered}
 y_0\in C_0,\qquad
 z_0\in\bar{\mathfrak n}\mathfrak cC_0I/\bar I,\\
 y_0+z_0b_p\,\bar i_p/i_p\equiv0\pmod{\mathfrak p}
                    \qquad(p\mid D),
 \end{gathered}
\end{equation}
where \(i_p\) generates \(I\) locally.
To check this, away from \(q\) test
\[
 u_0t+vb_s(\bar i_s/i_s)\bar t\in C_{0,s}
                    \qquad(t\in O_s).
\]
At étale places the linear and antilinear conditions separate.
The split level condition on \(\mathfrak n^{-1}O_s/O_s\)
adds the factor \(\bar{\mathfrak n}\) on the antilinear side.
At an odd ramified prime, set \(v_s^*=vb_s\bar i_s/i_s\).
Testing on \(1,\sqrt{-D}\) requires \(u_0+v_s^*\) and
\(\sqrt{-D}(u_0-v_s^*)\) to be integral; these are the two
inverse-different conditions and the extra congruence in
\eqref{tame:quaternion-lattice}. Here \(C_{0,s}=O_s\).
At \(q\), integrality separates because the linear and
antilinear norm valuations have opposite parities.

The degree identity is
\[
 Dma_C=\mathrm Ny_0+c_{\rm B}\mathrm Nz_0.
\]
Set
\[
 Z=(z_0)/(\bar{\mathfrak n}\mathfrak cC_0I/\bar I),
 \qquad b=q\,\mathrm NZ.
\]
Because \(\mathrm N\mathfrak c=q/c_{\rm B}\), this gives
\(a+Nb=mD\), with \(a=\mathrm Ny_0/a_C\), exactly as in the
coefficient formula. Neither component vanishes: a purely
linear map violates the inert-\(\ell\) degree condition,
whereas a purely antilinear map would force \(N\mid mD\).

\emph{Genus conditions and the number of maps.}
Fix \(y_0,C_0\). For an ideal \(Z\) of norm \(b/q\), the
condition that a generator \(z_0\) exist is that
\([Z\bar{\mathfrak n}\mathfrak cC_0]\) be a square.
Its norm is \(bNa_C/c_{\rm B}\). At \(p\mid D\), the quaternion
algebra splits, so \(-c_{\rm B}\) is a local norm.
After moving the ideal away from \(D\) by principal scaling,
Gauss genus theory identifies the square condition with
\[
 \varepsilon_{K,p}(-bNa_C)=1\qquad(p\mid D).
\]
If soluble, there are \(2^{\#\{p\mid D\}-1}\) classes \(I\), and
two generators \(z_0\) for each before the congruences.
At \(p\mid b\) the congruence in
\eqref{tame:quaternion-lattice} is automatic. At the other
discriminant primes the nonzero residues match up to a sign.

These signs are jointly equidistributed. Replacing \(I\) by
\(\mathfrak pI\) leaves \(I/\bar I\) unchanged and flips the
generator-ratio sign only at \(p\). Principal changes of
representative transport \(z_0\) by the corresponding quotient
with its conjugate and preserve the test. The relation obtained
by multiplying all ramified ideals is the principal ideal
\((\sqrt{-D})\); its transport changes the generator by
\(\sqrt{-D}/\overline{\sqrt{-D}}=-1\).
Thus that relation is absorbed by the two generator choices.
Of the \(2^{\#\{p\mid D\}}\) pairs, the
\(\#\{p\mid D:p\nmid b\}\) nonautomatic signs leave
\(2^{\#\{p\mid D:p\mid b\}}\) possibilities.
Summing over ideals \(Z\) gives
\[
 2^{\#\{p\mid D:p\mid b\}}\rho(b/q)
\]
homomorphisms per summand. Moreover \(d_q\) is odd, since
\(b=q\,\mathrm NZ\) and ideal norms have even valuation at
an inert prime.

\emph{The length of a lift.}
Each counted homomorphism has lifting length \((1+d_q)/2\).
Use the Grothendieck--Messing filtration criterion and
Serre--Tate theory over \(\mathcal W/q^j\), with the canonical
nilpotent divided powers on \(q\); here \(q\ge5\).
The integral first crystalline cohomology splits into two
unramified CM lines. Up to units, Frobenius has
\[
 F(e_1)=e_2,\qquad F(e_2)=qe_1;
\]
the Hodge summand is \(e_2\). Source and target have the same
CM type. The antilinear component interchanges the lines,
say \(e_1\mapsto Ae_2\), \(e_2\mapsto Be_1\).
Commutation with Frobenius gives \(v_q(B)=v_q(A)+1\).
Their sum is \(d_q\), the determinant valuation of this
component, so preservation of the Hodge line modulo \(q^j\)
is exactly
\[
 j\le v_q(B)=\frac{1+d_q}{2}.
\]
Prime-to-\(q\) level adds no lifting obstruction. This is also
the inert lifting exponent in
\cite[Corollary~7.15]{Conrad03}.
Combining the map count, this length, and the residue-degree
cancellation gives the negative of \eqref{eq:T6}, with its
stated normalization.
\end{proof}

\subsection{Geometric comparison and the logarithmic formula}

\begin{proof}[Proof of Proposition~\ref{tame:spectral}]
By Proposition~\ref{tame:intersection}, \(B_h(n)\) computes,
up to the common sign, the partial character-weighted
intersection away from \(r\). We explain why the combinations
in \eqref{eq:T4} compute the full heights. In particular, we
do not impose an unramified-character description downstairs
at \(r\).

The difference \(C_h(n)=B_h(n)-B_{h-2}(n)\) corresponds to
\(T_n(T_{r^h}-T_{r^{h-2}})\) on the second divisor. Only cyclic
\(r^h\)-paths remain: a noncyclic kernel contains full \(r\)-torsion
and factors through multiplication by \(r\). Exactly two cyclic
kernels at each full CM object remain at maximal order, one
on each split factor. Their Hilbert sums are both the original
sum, independently of \(h\ge1\), including the corresponding
cusp degrees. The second difference on the left of
\eqref{eq:T4} removes these terms.

Every remaining Heegner orbit before \(T_n\) has conductor
a positive power \(r^j\). Indeed a cyclic over-lattice that
remained maximal locally would be \(O_r\)-stable and lie on
one split factor. Ring class theory therefore places its field
of definition over \(K\) over \(H_{r^j}\). Write the orbit
with its multiplicity as a trace from such a field \(F'\),
including its own cusp subtraction, and pull the first divisor
up by the projection formula. The \(r\)-intersections of these
pulled-up terms are zero. Their reductions are ordinary,
since \(r\) is split, and both objects before \(T_n\) have full
order at \(\ell\). A matching homomorphism of degree \(n\)
would violate the same linear or conjugate-semilinear norm
obstruction at the inert prime \(\ell\). The cusps also
contribute nothing. We may consequently add the \(r\)-place
intersections as zero. The character over \(F'\) is now
unramified, so Lemma~\ref{tame:heights} identifies the sum with
\(h_i^{F'}\), using the normalized norm valuations in the
projection formula.

Applying \(\mathcal A\) to \eqref{eq:T4} at \(n=\ell\) inserts
\(\mathcal A T_\ell\) between the height arguments. The
prime-to-\(Nr\) Hecke recurrences in \(n\) commute with the
cyclic-path expansion already made. On the Jacobian,
\[
 \mathcal A T_\ell
     =\frac{\mathcal A_fa_f(\ell)}{\delta}\phi^*\phi_*
\]
rationally: \(\phi^*\phi_*/\delta\) is the orthogonal projector
onto the elliptic factor. Biextension functoriality over each
\(F'\) projects the pairing to \(E\), with this scalar;
only fixed denominators are cleared. Multiply the first
argument by the fixed \(2^b\) of Lemma~\ref{tame:heights},
and retrace the second to \(K\).

The resulting second arguments on \(E\) are the indicated
Hecke second differences of \(Q_X\), up to bounded-order cusp
torsion. Each cyclic operator has eigenvalue
\(a_f(r^h)-a_f(r^{h-2})\). After passage to the limit and the
additional difference \(h\mapsto h+2\), the geometric side is
\begin{equation}\label{tame:geometric-side}
 \pm\mathcal A_fa_f(\ell)(a_*-2)
       (\gamma_{h+2}-\gamma_h)\,
       \frac{h_\lambda^K(Q_X,Q_X)}{\delta}.
\end{equation}
Here the second difference uses
\(\gamma_{h+1}-2\gamma_h+\gamma_{h-1}=(a_*-2)\gamma_h\).
Cusp torsion disappears in rational heights; the fixed
multipliers cancel in these rational equalities.

On the other side, apply \eqref{eq:T5} at \(h+1,h,h-1\),
with coefficients \(1,2,1\), as required by \eqref{eq:T4}.
The determinant-one recurrence now inserts the factor
\(a_*+2\). If \(\alpha,\alpha^{-1}\) are the limiting
eigenvalues, then \(\gamma_h=\alpha^h+\alpha^{-h}\).
Their distinct non-root-of-unity property ensures that
\(\gamma_{h+2}-\gamma_h\ne0\) for a sufficiently large \(h\).
Cancel this quantity and \(\mathcal A_fa_f(\ell)\ne0\).
The remaining scalar
\[
 \frac{(a_*+2)(2-a_*)}{(a_*-2)(a_*+2)}=-1
\]
on the holomorphic side gives exactly
\eqref{tame:spectral-identity}. If \(Q_X\) is torsion its
rational height is zero, while
\(C\Omega_0L(f\otimes\varepsilon,1)\ne0\); hence \(Z'=0\).
No rational spanning assertion has entered this proof.
\end{proof}

\begin{remark}\label{tame:unknown-rank-direction}
The spectral comparison itself requires neither a nonzero cohomology
class nor the prescription \(\eta\ne0\). Omitting these data from the
open-image and Chebotarev construction still gives all the spectral
hypotheses of Definition~\ref{tame:data}. One may also impose extra
fixed cochain data, as in the unknown-center argument below.
Proposition~\ref{tame:spectral}, and in particular its
torsion implication, applies to every such direction.
This is the form used at an initially unknown analytic center.
\end{remark}

\begin{proof}[Proof of Theorem~\ref{tame:logarithm}]
Now assume \(\an(E)=1\). The companion has analytic rank zero,
so the classical low-rank theorem gives the rational rank
decomposition over \(K\). In particular \(Q_X\) is a rational
multiple of \(P\) modulo torsion; write \(Q_X=t_P P\) in
\(E(K)\otimes\Q\). The exact split Gross--Zagier normalization
\eqref{eq:GZ}--\eqref{eq:GZ-E} gives
\begin{equation}\label{tame:classical-height}
 C L'(f,1)L(f\otimes\varepsilon,1)
       =\frac{2H_{\rm abs}(Q_X)}{\delta}
       =\frac{2t_P^2H(P)}{\delta}.
\end{equation}
This is the single Hilbert-class trace and the absolute
Poincaré-height diagonal, with the restricted polarization;
see \cite[Chapters~I and V]{GZ86} and \cite{CST14}.
Together with \eqref{tame:period-rationality}, it also proves
the rationality of \(L'(f,1)/(\Omega_0H(P))\).

Choose the direction of Lemma~\ref{tame:direction}, so that
\(h_\lambda^\Q(P,P)\ne0\). By Lemma~\ref{tame:heights},
\[
 h_\lambda^K(Q_X,Q_X)=2t_P^2h_\lambda^\Q(P,P).
\]
Substitution of \eqref{tame:classical-height} into
\eqref{tame:spectral-identity} therefore gives
\[
 Z'=\pm d_0(2-a_*)\,
       \frac{L'(f,1)}{\Omega_0H(P)}
          h_\lambda^\Q(P,P).
 \tag{T3}\label{eq:T3}
\]
The factor \(2\) from field extension is the same factor as
in \eqref{tame:classical-height}; neither is discarded.
Finally combine \eqref{eq:T3} with the cup reciprocity formula
\eqref{eq:T2}. Since \(2-a_*\), \(h_\lambda^\Q(P,P)\), and
\(\log_\omega P\) are nonzero, cancellation yields
\eqref{eq:T1}.
\end{proof}
